% concatenated from 2ket_2.tex
%%%%%%%%%%%%%%%%%%%%%%%%%%%%%%%%%%%%%%%%%%%%%%%%%%%%%%%%%%%%%%%%%%%%%%%%
%%%%%%%%%%%%%%%%%%%%%%%%%%%%%%%%%%%%%%%%%%%%%%%%%%%%%%%%%%%%%%%%%%%%%%%%
%%  gmcfnf.tex
%%  LaTeX2e-file
%%  version of Jan 14, 2013
%%%%%%%%%%%%%%%%%%%%%%%%%%%%%%%%%%%%%%%%%%%%%%%%%%%%%%%%%%%%%%%%%%%%%%%%
%%%%%%%%%%%%%%%%%%%%%%%%%%%%%%%%%%%%%%%%%%%%%%%%%%%%%%%%%%%%%%%%%%%%%%%%
\documentclass[a4paper,12pt]{article}
%\special{papersize=17cm,24cm}
\usepackage{amssymb}
\usepackage{amsthm}
\usepackage{amsmath}
\usepackage{graphicx}
\usepackage{psfrag}
\usepackage{array,hhline}
%\usepackage[notref,notcite]{showkeys}
\usepackage{enumitem}
\usepackage{caption}
\usepackage{subcaption}

%%%%%%%%%%%%%%%%%%%%%%%%%%%%%%%%%%%%%%%%%%%%%%%%%%%%%%%%%%%%%%%%%%%%%%%%
%%  Theorems and Definitions

\newtheorem{lemma}{Lemma}[section]
\newtheorem{theorem}[lemma]{Theorem}
\newtheorem{proposition}[lemma]{Proposition}
\newtheorem{conjecture}[lemma]{Conjecture}
\newtheorem{corollary}[lemma]{Corollary}
\theoremstyle{definition}
\newtheorem{definition}[lemma]{Definition}
\newtheorem{remark}[lemma]{Remark}
\newtheorem{example}[lemma]{Example}

\numberwithin{equation}{section}
\numberwithin{figure}{section}

\renewcommand{\theenumi}{\roman{enumi}}
\renewcommand{\labelenumi}{{\rm(\theenumi)}}

\newcommand{\df}[1]{{\em #1}}
\newcommand{\wti}[1]{{\widetilde #1}}
\newcommand{\NN}{{\mathbb N}}
\newcommand{\bsx}{{\boldsymbol x}}
\newcommand{\bsy}{{\boldsymbol y}}
\newcommand{\bsz}{{\boldsymbol z}}
\newcommand{\nix}{{\mbox{\vphantom x}}}

\newcommand{\Aset}{\mathcal{A}}
\newcommand{\Bset}{\mathcal{B}}
\newcommand{\Cset}{\mathcal{C}}
\newcommand{\Dset}{\mathcal{D}}
\newcommand{\Eset}{\mathcal{E}}
\newcommand{\Fset}{\mathcal{F}}
\newcommand{\Gset}{\mathcal{G}}
\newcommand{\Lset}{\mathcal{L}}
\newcommand{\Mset}{\mathcal{M}}
\newcommand{\Nset}{\mathcal{N}}
\newcommand{\Pset}{\mathcal{P}}
\newcommand{\Sset}{\mathcal{S}}
\newcommand{\Tset}{\mathcal{T}}
\newcommand{\Xset}{\mathcal{X}}
\newcommand{\Yset}{\mathcal{Y}}

\newcommand{\ie}{\emph{i.e.}}
\newcommand{\eg}{\emph{e.g.}}


%%%%%%%%%%%%%%%%%%%%%%%%%%%%%%%%%%%%%%%%%%%%%%%%%%%%%%%%%%%%%%%%%%%%%%%%
%%%%%%%%%%%%%%%%%%%%%%%%%%%%%%%%%%%%%%%%%%%%%%%%%%%%%%%%%%%%%%%%%%%%%%%%
\begin{document}

%%%%%%%%%%%%%%%%%%%%%%%%%%%%%%%%%%%%%%%%%%%%%%%%%%%%%%%%%%%%%%%%%%%%%%%%
\title{\large{On the usage of $2$-node lines in $n$-correct and $GC_n$ sets}}

\author{H. Hakopian, \  G. Vardanyan, \  N. Vardanyan}

\author{H. Hakopian, G. Vardanyan, N. Vardanyan\\ \\
 \textit{Yerevan State University, Yerevan, Armenia}\\   \textit{Institute of Mathematics of NAS RA}}
\date{}
\maketitle

%%%%%%%%%%%%%%%%%%%%%%%%%%%%%%%%%%%%%%%%%%%%%%%%%%%%%%%%%%%%%%%%%%%%%%%%
\begin{abstract} 
An n-correct set $\Xset$ in the plane is a set of nodes admitting unique
 interpolation with bivariate polynomials of total degree at most $n$.
A $k$-node line is a line passing through exactly $k$ nodes of $\Xset.$ A line can pass through at most $n+1$ nodes of an $n$-correct set. An $(n+1)$-node line is called maximal line (C. de Boor, 2007). We say that a node $A\in\Xset$ uses a line $\ell,$ if $\ell$ is a factor of the fundamental polynomial of the node $A.$

Let $\Xset$ be an $n$-correct set.
One of the main problems we study in this paper is to determine the maximum possible number of used $2$-node lines that share a common node $B \in \mathcal{X}$.

 
We show that this number equals $n$. Moreover,  if there are $n$ such $2$-node lines, then $\mathcal{X}$ contains exactly $n$ maximal lines not passing through the common node $B$. Furthermore, if $\mathcal{X}$ is $GC_n$ set, there exists an additional maximal line passing through $B$. Hence, in this case, $\mathcal{X}$ has $n+1$ maximal lines and is Carnicer–Gasca set of degree $n$. Note that Carnicer–Gasca sets of degree $n$ with a prescribed set of $n$ used $2$-node lines can be readily constructed.
\end{abstract}

{\bf Key words:}
Bivariate interpolation, Carnicer-Gasca set, $n$-correct set, $GC_n$ set, maximal line, maximal curve.


{\bf Mathematics Subject Classification 2020:} 41A05, 41A63, 14H50

%%%%%%%%%%%%%%%%%%%%%%%%%%%%%%%%%%%%%%%%%%%%%%%%%%%%%%%%%%%%%%%%%%%%%%%%
%%%%%%%%%%%%%%%%%%%%%%%%%%%%%%%%%%%%%%%%%%%%%%%%%%%%%%%%%%%%%%%%%%%%%%%%

\section{Introduction}
Denote by $\Pi_n$ the space of bivariate polynomials of total degree
$\le n.$ We set
$$
N:=\dim \Pi_n,\ \hbox{where}\ N=(1/2){(n+1)(n+2)}.
$$
\noindent Consider a set of $s$ distinct nodes in the plane:
$$\Xset:=\Xset_s=\{ (x_1, y_1), \dots , (x_s, y_s)\}.$$

The problem of finding a polynomial $p \in \Pi_n$ satisfying the conditions
\begin{equation}\label{int cond}
p(x_i, y_i) = c_i, \ \ \quad i = 1, 2, \dots s  ,
\end{equation}
for a given data $\bar c:=\{c_1, \dots, c_s\}$ is called \emph{interpolation problem.}
\begin{definition}
A set of nodes $\Xset_s$ is called
$n$-\emph{correct} if for any data $\bar c$ there exists a
unique polynomial $p \in \Pi_n$, satisfying the conditions
\eqref{int cond}.
\end{definition}
A necessary condition of
$n$-correctness is: $\#\Xset_s=s = N.$

Denote by $p|_\Xset$ the restriction of $p$ on $\Xset.$
\begin{proposition} \label{correctii}
A set of nodes $\Xset$ with $\#\Xset=N$ is $n$-correct if and only if
$$p \in \Pi_n,\ p|_\Xset=0\implies p = 0.$$
\end{proposition}
A polynomial $p \in \Pi_n$ is called an $n$-\emph{fundamental polynomial}
for $ A \in \Xset$ if
$$ p|_{\Xset\setminus\{A\}}=0\ \hbox{and}\ p(A)=1.$$
We denote an
$n$-fundamental polynomial of $A \in\Xset$ by $p_A^\star=p_{A,\Xset}^\star.$

Let us mention that sometimes we call $n$-fundamental a polynomial $p\in \Pi_n$ satisfying the conditions
$$ p|_{\Xset\setminus\{A\}}= 0\ \hbox{and}\ p(A)\neq 0,$$
since it equals a non-zero constant times $p_A^\star.$

\begin{definition}
A set of nodes $\Xset_s$ is called $n$-\emph{independent} if each node
has $n$-fundamental polynomial. 
\end{definition}
Fundamental polynomials are linearly independent. Therefore a
necessary condition of $n$-independence is $\#\Xset_s=s \le N.$


A \emph{plane algebraic curve} of degree $n,\ n\ge 1,$ is the zero set of some bivariate polynomial $p$ of degree $n,$ i.e., the set $\left\{A : p(A)=0\right\}.$ To simplify notation, we shall use the same letter,  say $p$,
to denote the polynomial $p$ and the curve given by the equation $p(x,y)=0$.
In particular, by $\ell$ we denote a linear
polynomial from $\Pi_1$ and the line defined by the equation
$\ell(x, y)=0.$

Let us mention that in the expressions like ${\mathcal X \cap p},\ \Xset\setminus p,$ by the polynomial $p\in\Pi_k$ we mean its zero set. 
\begin{definition} Let $\Xset$ be an $n$-correct set.
We say that a node $A\in\Xset$ uses a curve $q$ of degree $k\le n,$ if $q$ is a component of the fundamental polynomial of the node $A:$ $$ p^\star_A=q r,\ \hbox{where}\  r\in\Pi_{n-k}.$$
\end{definition}

\subsection{$\bf{GC_n},$ sets, Chung-Yao and Carnicer-Gasca sets}

Let us now consider a particular type of $n$-correct sets (see \cite{HV}) that satisfy a so-called geometric characterization (GC) property introduced by K.C. Chung and T.H. Yao \cite{CY77}:
\begin{definition}[\cite{CY77}]
An $n$-correct set ${\mathcal X}$ is called \emph{$GC_n$ set }
 if  the
$n$-fundamental polynomial of each node $A\in{\mathcal X}$ is a
product of $n$  linear factors.
\end{definition}
Thus, each node of $GC_n$ set uses exactly $n$ lines.

The following proposition is well-known (see e.g. \cite{HJZ09b}
Prop. 1.3):
\begin{proposition}\label{prp:n+1ell}
Suppose that a polynomial $p \in
\Pi_n$ vanishes at $n+1$ points of a line $\ell.$ Then, we have that
$
p = \ell  q  ,\ \text{where} \ q\in\Pi_{n-1}.
$
\end{proposition}
We say, that a line $\ell$ is a $k$-\emph{node line} if it passes through exactly $k$ nodes of $\Xset.$
The above proposition implies that at most $n+1$ nodes of an $n$-independent set  can be collinear. An $(n+1)$-node line $\ell$ is called a \emph{maximal line} (C. de Boor, \cite{dB07}).

Denote by $\Mset(\Xset)$ the set of maximal lines of an $n$-correct set $\Xset.$



Below some basic properties of maximal lines are presented.

\begin{proposition}[\cite{CG01}, Prop.2.1] \label{hatk} Let $\Xset$ be an $n$-correct set. Then:
\begin{enumerate}
\item
Any two maximal lines intersect at a node of $\Xset.$
\item
Any three maximal lines are not concurrent.
\item
$\#\Mset(\Xset)\le n+2.$
\end{enumerate}
\end{proposition}

 
Let a set $\Lset=\{\ell_0,\ldots,\ell_{n+1}\}$ of $n+2$ lines be in general position, that is 
no two lines of $\Lset$ are parallel, and 
no three lines of $\Lset$ are concurrent.

Then the set $\Xset$ of ${n+2\choose 2}$ intersection points
of these lines is called Chung-Yao set of degree $n$ (see \cite{HV}):
\begin{equation}\label{Aij}\Xset= \{A_{ij}: 0\le i<j\le n+1\},\ \hbox{where}\ A_{ij}=\ell_i\cap\ell_j.\end{equation} 
 
Note that each node $A_{ij}$ belongs to $2$ lines: $\ell_i$ and $\ell_j,$ and the product of the remaining $n$ lines gives $p^\star_{A_{ij}}.$ Thus $\Xset$ is $GC_n$ set. It is easily seen that  the lines $\ell_i,\ i=0,\ldots,n+1,$ are maximal for $\Xset.$ In view of Proposition \ref{hatk}, (iii), no other line is maximal for the set $\Xset,$ that is, $\Lset=\Mset(\Xset).$

On the other hand if $\Xset$ is an $n$-correct set and $\#\Mset(\Xset)= n+2$ then $\Xset$ is Chung-Yao set of degree $n.$ Indeed, assume that $\Mset(\Xset)= \{\ell_0,\ldots,\ell_{n+1}\}$ is the set of maximal lines of $\Xset.$ Then, by using Proposition \ref{hatk}, (i),(ii), we readily get \eqref{Aij}.

Thus 
 an $n$-correct set $\Xset$ is Chung-Yao set of degree $n$ if and only if $\#\Mset(\Xset)= n+2.$

Now, let a set $\Lset=\{\ell_0,\ldots,\ell_n\}$ of $n+1$ lines be in general position (see Fig.\ref{Fig1}, for $n=3$).
Then the set of ${n+1\choose 2}$ intersection points $\{A_{ij}: 0\le i<j\le n\}$
of these lines together with a set of $n+1$ other noncollinear points $\{A_i\in\ell_i: 0\le i\le n\},$ is called Carnicer-Gasca set of degree $n$ (see \cite{HV}, Section 3). 
\begin{figure}
\begin{center}
\includegraphics[width=8.0cm,height=5.cm]{Figure1.png}
\end{center}
\caption{Carnicer-Gasca sed of degree $3$} \label{Fig1}
\end{figure}


Denote by $\ell_{AB}$ the line between the points $A,B.$

Each node  $A_i$ here belongs to only one line: $\ell_i,$ and the product of the remaining $n$ lines gives $p^\star_{A_i}.$ Next, each node $A_{ij}$ belongs to two lines: $\ell_i$ and $\ell_j,$ and the product of the remaining $n-1$ lines and the line $\ell_{A_i,A_j}$ gives $p^\star_{A_{ij}}.$ Hence $\Xset$ is $GC_n$ set. Note that the lines $\ell_i,\ i=0,\ldots,n,$ are maximal for $\Xset.$  Thus, in Carnicer-Gasca set except the maximal lines  also the lines $\ell_{A_i,A_j}$ are used. 

It can be checked readily that  
 an $n$-correct set $\Xset$ is Carnicer-Gasca set of degree $n$ if and only if $\#\Mset(\Xset)= n+1.$








\subsection{Maximal curves}

Set $$d(n, k) := N_n - N_{n-k} =(n+1)+n+(n-1)+\cdots+(n-k+2)= (1/2) k(2n+3-k).$$

The following is a  generalization of Proposition \ref{prp:n+1ell}.
\begin{proposition}[\cite{Raf}, Prop. 3.1]\label{maxcurve}
Let $\Xset$ be an $n$-correct set and $q$ be an algebraic curve of degree $k \le n$ with no multiple components. Then, the following hold:

$(i)$ Any subset of $q$ containing more than $d(n,k)$ nodes of $\Xset$ is
$n$-dependent.

$(ii)$ Any subset of $q$ containing exactly $d(n,k)$ nodes of $\Xset$ is $n$-independent if and only if
\vspace{-.5cm}

$$\quad p\in {\Pi_{n}}\ \text{and}\ \ p|_{{\mathcal X}} = 0 \implies  p = qr,\ \hbox{where}\ r \in \Pi_{n-k}.$$
\end{proposition}
\noindent Thus at most $d(n,k)$ $n$-independent nodes lie in a curve $q$ of degree $k \le n$.
\begin{definition}\label{def:maximal}
Let $\Xset$ be an $n$-correct set of nodes. A curve of degree $k \le n$ passing through $d(n,k)$ points
of $\mathcal X$ is called maximal curve for $\Xset.$
\end{definition}
Clearly, a maximal curve cannot have a multiple component.

The following proposition gives a characterization of the maximal curves:
\begin{proposition}[\cite{Raf}, Prop. 3.3] \label{maxcor}
Let $\Xset$ be an $n$-correct set. Then
a curve $f\in\Pi$ of degree $k,\ k\le n,$ is a maximal curve for $\Xset$ if and only if 
it is used by any node of the set $\Xset\setminus f.$
\end{proposition} 

Note that this, in view of the Lagrange formula, follows directly from Proposition \ref{maxcurve}. 

One readily gets from here that for an $n$-correct set $\Xset:$ 
\begin{equation}\label{maxmax}q\in\Pi_k \hbox{ is a maximal curve}  \iff \Xset\setminus q  \hbox{ is an $(n-k)$-correct set.}\end{equation} 
One also obtains readily for  a $GC_n$ set $\Xset:$ 

If $f$ is a maximal curve of degree $k\le n,$ then $f$ is a product of $k$ distinct lines (see \cite{HVV}, Prop. 5.1). 

Next result concerns components of a maximal curve. We will use it in the special case of the line-components.
\begin{proposition}\label{maxcomp} Let $\Xset$ be an $n$-correct set  and
a curve $q$ of degree $k\le n,$ be a maximal curve. Suppose that a curve $s$ of degree $m\le k,$ is a component of $q,$ that is $q=s r, \deg r=k-m.$
Then $s$ passes through  at least $d(n-k+m,m)$ nodes of $\Xset \setminus r.$ Moreover, $s$ passes through  exactly $d(n-k+m,m)$ nodes of $\Xset \setminus r$ if and only if $r$ is a maximal curve of degree $k-m.$ 
\end{proposition}

\begin{proof} Assume, by way of contradiction, that   $s$ passes through  less than $d(n-k+m,m)$ nodes of $\Xset \setminus r.$
Then, since $q=s r$ passes through $d(n,k)$ nodes of $\Xset,$ we get that the curve $r, \deg r=k-m,$ passes through more than $d(n,k)-d(n-k+m,m)=d(n,k-m)$ nodes of $\Xset,$ which is contradiction.

Now, if $s$ passes through  exactly $d(n-k+m,m)$ nodes of $\Xset \setminus r$ then above calculation shows that  $r$ passes through exactly $d(n,k)-d(n-k+m,m)=d(n,k-m)$ nodes of $\Xset.$ Hence, $r$ is a maximal curve of degree $k-m.$ And, vice versa, if $r$ passes through exactly $d(n,k-m)$ nodes of $\Xset$ then we similarly get that  $s$ passes through  exactly $d(n-k+m,m)$ nodes of $\Xset \setminus r.$  
\end{proof}

In the sequel we will need the following
\begin{theorem} [\cite{HakTor}, Thm. 4.1] \label{uncurve} Assume that $\Xset$ is an n-independent set of
$d(n; k-1)+2$ nodes lying in a curve of degree $k\le n.$ Then the curve
is determined uniquely by these nodes.
\end{theorem}

\subsection{Special triplets}

Let us start with the following
\begin{definition}
Let $\Xset$ be an $n$-correct set of nodes. A triple of nodes $\{A,B,C\}\subset\Xset$ is called a \emph{special triplet} if there exists a maximal curve $f$ of degree $n-1$ such that 
\begin{equation} \label{020} \Xset\setminus f =\{A,B,C\}.\end{equation}
\end{definition}
We call above $f$ the maximal curve associated with $\{A,B,C\}.$

An example of a special triplet in a $GC_5$ set  is presented in Fig.\ref{Fig2}. The associated maximal curve in this case is:
$f(x,y)=xy(y-1)(x+y-5).$


An example of a nonspecial triplet in a $GC_5$ set  is presented in Fig.\ref{Fig3}. Indeed, assume by way of contradiction, that the triplet is special and $f$ is the associated maximal curve. Then, in view of Proposition \ref{prp:n+1ell}, we get that $f(x,y)=xy(x+y-5)\ell(x,y),$ where $\ell\in\Pi_1.$ Clearly we obtain also that $\ell$ vanishes at $(1,2), (2,1)$ and $(2,2).$ Therefore $\ell=0$ and $f=0,$ which is  contradiction. 

\begin{figure}
     \centering
     \begin{subfigure}[b]{0.3\textwidth}
         \centering
         \includegraphics[width=\textwidth]{Figure2.png}
         \caption{A special triplet}
         \label{Fig2}
     \end{subfigure}
     \hfill
     \begin{subfigure}[b]{0.3\textwidth}
         \centering
         \includegraphics[width=\textwidth]{Figure3.png}
         \caption{A nonspecial tripliet}
         \label{Fig3}
     \end{subfigure}
          \caption{Two triplets in a $GC_5$ set}
        \label{fig2}
\end{figure}


\begin{proposition} \label{sptr} Let $\Xset$ be an $n$-correct set of nodes and $\{A,B,C\}\subset\Xset$ be a special triplet. Then the nodes $A,\ B,$ and $C$ are noncollinear and use the lines $\ell_{BC},\ \ell_{CA},$ and $\ell_{AB},$ respectively.
\end{proposition}
\begin{proof} Let $f$ be the associated with $\{A,B,C\}$ maximal curve. Assume, by way of contradiction, that $A,B,C$ are collinear, i.e., belong to some line $\ell.$ Then the polynomial $\ell f\in\Pi_n$ vanishes on $\Xset,$ which contradicts Proposition \ref{correctii}. Next, we readily obtain that
\begin{equation}\label{ABC} p^\star_{A,\Xset}=\ell_{BC} f, \ \ p^\star_{B,\Xset}=\ell_{CA} f,\ \ \hbox{and} \ \  p^\star_{C,\Xset}=\ell_{AB} f.\end{equation}
\end{proof}
Note that the triplet in Fig.\ref{Fig2} satisfies the condition \eqref{ABC}, though it is not a special triplet.

\begin{proposition} \label{sptr2} Let $\Xset$ be an $n$-correct set of nodes. Let also $\{A,B,C_1\}$ and $\{A,B,C_2\}$ be special triplets, where $A,B,C_1,C_2 \in\Xset$. Then $C_1=C_2.$
\end{proposition}
\begin{proof} Suppose that $f_i\in\Pi_{n-1}$ be the maximal curve associated with $\{A,B,C_i\},$  $\ i=1,2.$ Then $f_1$ and $f_2$ pass through $N-4$ nodes of the set
$\Xset\setminus \{A,B,C_1,C_2\}.$ From here, in view of Theorem \ref{uncurve}, we conclude that $f_1= f_2$ and therefore $C_1=C_2,$ since $d(n,n-2)+2=N-4.$

\end{proof}
\begin{lemma}\label{cor:2ket}
Let $\Xset$ be an $n$-correct set and $\ell_{AB}$ be a $2$-node line used by a node $C,$ where $A, B,C\in\Xset.$ Then $\{A,B,C\}$ is a special triplet. 
In this case, we say that the special triplet $\{A, B, C\}$ arises from the 2-node line $\ell_{AB}.$
\end{lemma}
\begin{proof}
The node $C$ uses the  line $\ell_{AB}$ means that
\begin{equation}\label{01}p^\star_{C,\Xset}=\ell_{AB} f,\ \hbox{where}\  f\in\Pi_{n-1}.
\end{equation}
Since $\ell_{AB}$ is a $2$-node line we conclude from here that the relation \eqref{020} takes place. Therefore $\{A,B,C\}$ is a special triplet and $f$ is the associated maximal curve of degree $n-1$.
\end{proof}
Note that the triplet in Fig.\ref{Fig2} is not arisen from $2$-node line. Indeed, the lines through the vertices of the triplet are $5,5$ and $4$-node lines.

Next, we obtain the following result of Carnicer and Gasca:
\begin{corollary}[\cite{CG03}, Prop. 4.2] \label{2nodeuse} Let $\Xset$ be an $n$-correct set. Then any $2$-node line can be used by at most one node.
\end{corollary}
Indeed, assume that a $2$ node line $\ell_{AB}$ is used by two nodes $C_1,C_2\in\Xset.$ Then, by Lemma \ref{cor:2ket}, we get that $\{A,B,C_1\}$ and $\{A,B,C_2\}$ are special triplets. Now, by  Proposition \ref{sptr2}, we obtain that $C_1=C_2.$ 



Let $\Xset$ be an $n$-correct set.
One of the main problems we study in this paper is to determine the maximum possible number of used $2$-node lines that share a common node $B \in \mathcal{X}$.


 
We show that this number equals $n$. Furthermore, if there are exactly $n$ such $2$-node lines, then $\mathcal{X}$ contains precisely $n$ maximal lines not passing through the common node $B$. Moreover, if $\mathcal{X}$ is $GC_n$ set, then there exists an additional maximal line passing through $B$. Hence, in this case, $\mathcal{X}$ has $n+1$ maximal lines and is Carnicer–Gasca set of degree $n$. Finally, note that Carnicer–Gasca sets of degree $n$ with a prescribed set of $n$ used $2$-node lines can be readily constructed.

Indeed, suppose $\mathcal{X}$ is Carnicer–Gasca set of degree $n$ with  maximal lines $\ell_0,\ldots,\ell_n.$ Thus the intersection nodes $A_{ij}$  are determined as described earlier. Then let us choose the other nodes $A_i.$ Consider a node $B := A_0 \in \ell_0.$ Choose the remaining nodes $A_i \in \ell_i,\ i=1,\ldots,n,$ so that the $n$ lines $\ell_{A_iB}$ are $2$-node lines (see Fig. \ref{Fig1} for $n=3$). 
By the earlier observation, these lines are used by the nodes $C_i:=A_{0i},\ i=1,\ldots,n,$ respectively. Thus we have exactly $n$ used $2$-node lines that share the node $B\in\Xset.$ Note that the maximal lines not passing through the node $B$ are the lines $\ell_i\equiv\ell_{A_iC_i},\ i=1,\ldots,n.$ And the $(n+1)$-th maximal line, passing through $B,$ is the line $\ell_0\equiv\ell_{BC_1}$ to whom ($B$ and) the above nodes $C_i$ belong. 










\section{Main results}
Now, we are in a position to present a main result.
\begin{theorem}\label{th1sptr}
Let $\Xset$ be an $n$-correct set and $\mu:=\mu_k$ be a maximal curve of degree $k.$ Hence the set $\Bset:=\Bset_{n-k}=\Xset\setminus \mu$ is $(n-k)$-correct set. Denote by $\Aset:=\Aset_{k}=\Xset\cap \mu.$\\ Then a triplet of nodes $\{A,B,C\},$ where $A \in \Aset,\ B\in \Bset,$ and $C \in \Xset,$ is a special triplet if and only if the following two statements take place.
\begin{enumerate}
\setlength{\itemsep}{0mm}
\item
The line $\ell_1:=\ell_{AC}$ is a component of $\mu:\ \mu=\ell_1\mu_{k-1},$ in particular $B\notin \ell_1$ and $C\in \Aset.$  
\item
The set $\dot{\bar \ell}_1:=\{M\in{\ell_1}\cap  \Xset: \mu_{k-1}(M)\neq 0\}$ besides the nodes $A$ and $C,$ contains exactly $n-k$ more nodes of $\Xset,$ which 
coincide with the zeroes of $p^\star_{B,\Bset}\in\Pi_{n-k}$ on the line $\ell_1.$ Hence the latter zeroes are all real and distinct.
\end{enumerate}
Moreover, the statements \emph{(i) and (ii)} imply the following:
\begin{enumerate}
\setcounter{enumi}{2}
\setlength{\itemsep}{0mm}
\item
The curve $\mu_{k-1}:=\frac{1}{\ell_1}\ \mu$ is a maximal curve of degree $k-1.$ 

\item
The set $\Bset_{n-k+1}:=\Xset\setminus\mu_{k-1}=\Bset\cup \dot{\bar \ell}_1$ is an  $(n-k+1)$-correct set, for which $\ell_1$ is a maximal line.  
\item
$\Aset_{k-1}:=\Xset\cap\mu_{k-1}=\Aset\setminus \dot{\bar \ell}_1.$
\item 
$\{A,B,C\}$ is a special triplet in $\Bset_{n-k+1}$ too.
 
\item
If $\{B,D,E\}\subset\Xset$  is a special triplet, where $D\in \dot{\bar \ell}_1,$ then $\{B,D,E\}=\{A,B,C\},$ i.e., $\{D,E\}=\{A,C\}.$
\end{enumerate}
\end{theorem}


\begin{proof}
We have that
\begin{equation}\label{222}p^\star_{B,\Xset}=p^\star_{B,\Bset}\ \mu.
\end{equation}
On the other hand, Lemma \ref{sptr} implies that
\begin{equation}\label{202}p^\star_{B,\Xset}=\ell_1 \tilde\mu,
\end{equation}
where $\tilde\mu:=\tilde\mu_{n-1}$ is the  maximal curve of degree $n-1$ associated with $\{A,B,C\}:$
\begin{equation}\label{0207}\Xset\setminus \tilde\mu=\{A,B,C\}.\end{equation}
In view of \eqref{222} and \eqref{202} the line $\ell_1$ is either a component of $\mu$  or a component of  $p^\star_{B,\Bset}.$ In the latter case we readily conclude that $\mu$ divides $\tilde\mu,$ that is $\tilde\mu=q\mu,$ where
$q\in\Pi_{n-k-1}.$ Therefore, we obtain that 
$$\tilde\mu(A)=q(A)\mu(A)=0,$$  
which contradicts \eqref{0207}.



Thus $\ell_1=\ell_{AC}$ is a component of $\mu_k$ and (i) is proved.

Now denote $\mu_{k-1}:=\frac{1}{\ell_1}\ \mu_k,\ \mu_{k-1}\in\Pi_{k-1}.$ From the relations \eqref{222} and \eqref{202} we obtain that 
\begin{equation}\label{030}\tilde\mu=p^\star_{B,\Bset}\ \mu_{k-1}.\end{equation}


Now, let us consider the set of nodes  given in (ii):
$$\dot{\bar{\ell}}_1:=\{M\in{\ell_1}\cap  \Xset: \mu_{k-1}(M)\neq 0\}.$$


We have that ${\dot{\bar{\ell}}}_1\subset \ell_1\cap\Xset.$ Let us verify that $A,C\in \dot{\bar{\ell}}_1.$ 

Indeed, in view of \eqref{0207}, we have that $\tilde\mu$ does not vanish on $A$ and $C.$ Therefore  we obtain from \eqref{030} that $\mu_{k-1}$ does not vanish either.

Now,  \eqref{030} implies that besides $A$ and $C$ the set $\dot{\bar{\ell}}_1$ may contain only zeroes of 
the polynomial $p^\star_{B,\Bset}\in\Pi_{n-k}$ on the line $\ell_1,$ hence, at most $n-k$ nodes.

On the other hand the line $\ell_1$ is a component of the maximal curve $\mu_k$ of degree $k.$ Therefore, according to Proposition \ref{maxcomp}, with $m=1,$ besides $A$ and $C$ the set $\dot{\bar{\ell}}_1$ must have at least $d(n-k+m,m)-2=n-k+2-2=n-k$ nodes from $\Xset.$ 
Thus we conclude that the polynomial $p^\star_{B,\Bset}$ has exactly $n-k$ distinct real zeroes on the line $\ell,$ different from $A$ and $C,$ and all they belong to $\dot{\bar{\ell}}_1.$ Thus the statement (ii) is proved.

From here, in view of Proposition \ref{maxcomp}, with $m=1,$ we obtain that the curve $\mu_{k-1}$ is a maximal curve of degree $k-1.$ Then, the relation \eqref{maxmax} implies that $\Bset_{n-k+1}$ is an  $(n-k+1)$-correct set, for which thus $\ell_1$ is a maximal line. The remaining equalities in (iv) and (v) follow from the definitions of $\mu_{k-1}$ and ${\dot{\bar{\ell}}}_1.$  Hence, the statements (i) and (ii) imply (iii), (iv) and (v). 

To verify the statement (vi) it suffices to prove that $p^\star_{B,\Bset}\in\Pi_{n-k}$ is the 
associated maximal curve of the triplet $\{A,B,C\}:$ 
\begin{equation}\label{1.2.3}B_{n-k+1}\setminus p^\star_{B,\Bset}=\{A,B,C\}.
\end{equation}
Note that this readily follows from the statements (ii) and (iv).

Next, the statement (vii) follows from the following consequence of the statement (ii): \\ 
If $\{B,D,E\}\subset\Xset,$ where $D\in \dot{\bar \ell}_1,$ is a special triplet, then $p^\star_{B,\Bset}(D)\neq 0.$

Indeed, from here we obtain that either $D=A$ or $D=C.$ This, in view of Proposition \ref{sptr2}, implies that $E=C$ or $E=A,$ respectively.


Finally, let us assume that the statements (i), (ii) take place and prove that the triplet $\{A,B,C\}$ is special in $\Xset.$ Indeed, according to (ii), the curve  $\mu_{k-1}:=\frac{1}{\ell_1}\ \mu\in\Pi_{k-1}$ does not vanish at exactly $n-k+2$ nodes of the line $\ell_1$, namely at the nodes of the set $\dot{\bar \ell}_1.$   Now, (ii) and the relation \eqref{1.2.3} imply  that the polynomial 
$$p^\star_{B,\Bset}\ \mu_{k-1}\in \Pi_{n-1}$$
 vanishes at all the nodes of the set $\Xset\setminus\{A,B,C\}.$

Therefore  $\{A,B,C\}$ is a special triplet in $\Xset.$
\end{proof}


\subsection{The case of several special triplets}

Next, let us consider the case when there are two or more special triplets with a common node. 

\begin{theorem}\label{prp:m2ket}
Let $\Xset, \mu:=\mu_k, \Bset:=\Bset_{n-k},\Aset:=\Aset_{k},$ be as in Theorem \ref{th1sptr}.
Let $\Xset, \mu:=\mu_k, \Bset:=\Bset_{n-k},\Aset:=\Aset_{k},$ be as in Theorem \ref{th1sptr}. Suppose that the distinct triplets $\{A_i,B,C_i\},\  1\le i\le m,$ with a common node $B$ are special, where $m\ge 2, A_i \in \Aset, B\in \Bset,$ and $C_i \in \Xset.$ 

Then, the following statements take place for each $i,\ 1\le i\le m:$
\begin{enumerate}
\setlength{\itemsep}{0mm}
\item
The line $\ell_i:=\ell_{A_iC_i}$ is a component of $\mu_{k-i+1},$ in particular  $B\notin\ell_i$ and $C_i\in\Aset_{k-i+1}.$  These $m$ components are distinct, hence $m\le k.$ Moreover, we have that
$$\{A_j,C_j\}\cap\ell_{A_iC_i} =\emptyset\ \ \forall\ 1\le j\le m,\ j\neq i.$$
\item
The set $\dot{\bar\ell}_i:=\{M\in{\ell_i}\cap  \Xset: \mu_{k-i}(M)\neq 0\}$ besides the nodes $A_i$ and $C_i,$ contains exactly $n-k+i-1$ nodes of $\Xset,$ from which $n-k$ coincide with the zeroes of $p^\star_{B,\Bset}\in\Pi_{n-k}$ on the line $\ell_i,$ which are all real and distinct, and $i-1$ are the following points of intersection: 
$$D_{ij}:=\ell_i\cap\ell_j\ \ \forall\ 1\le j \le i-1.$$ 
\item
The curve $\mu_{k-i}:=\frac{1}{{\ell_i}}\mu_{k-i+1}=\frac{1}{\ell_1\cdots\ell_i}\ \mu$  is a maximal curve of degree $k-i.$  
\item
The set  
$\Bset_{n-k+i}:=\Xset\setminus\mu_{k-i}=\Bset_{n-k+i-1}\cup \dot{\bar\ell}_i=\Bset\cup\dot{\bar\ell}_1\cup\cdots\cup\dot{\bar\ell}_i
$ 

is an  $(n-k+i)$-correct set, for which $\ell_i$ is a maximal line.
\item
$\Aset_{k-i}:=\Xset\cap\mu_{k-i}=\Aset_{k-i-1}\setminus\dot{\bar \ell}_i=\Aset\setminus (\dot{\bar \ell}_1\cup\cdots\cup\dot{\bar \ell}_i).$ 

\item
The line $\ell_i$ is a maximal line in  $\Bset_{n-k+s}$, where $i\le s\le m.$ In particular all the lines $\ell_1,\ldots,\ell_m$ are maximal lines in $\Bset_{n-k+m}.$
 \end{enumerate}
\end{theorem}
\begin{proof}
Let us use indiction on $m$ for the first two statements of Theorem.
The case $m=1$ follows from Theorem \ref{th1sptr}.
Let us assume that statements (i) and (ii) of Theorem are valid for $m$ and prove them for $m+1.$
Later we will prove, as in the case of Theorem \ref{th1sptr}, that statements (i) and (ii) imply the remaining statements (iii)-(vi). 

Let $\{A_i,B,C_i\},$ be the distinct special triplets $1\le i\le m+1,$ where $A_i \in \Aset,\ B\in \Bset,$ and  $C_i \in \Xset.$ 
Since the triplets are distinct, Proposition \ref{sptr2} readily implies that 
\begin{equation*}\label{ACij}\{A_i,C_i\}\cap \{A_j,C_j\}=\emptyset\ \ \forall\ 1\le i<j\le m+1.\end{equation*}

Now, let us check that $A_{m+1}\in \Aset_{k-m}.$

Indeed, in view of  Theorem \ref{th1sptr}, (vii), with $\{A,B,C\}=\{A_i,B,C_i\},\ 1\le i\le m,$ and $\{B,D,E\}=\{A_{m+1},B,C_{m+1}\},$ we conclude that $A_{m+1}\notin \dot{\bar\ell}_i,$ since $A_{m+1}$ differs from $A_i$ and $C_i.$  Hence (v), with $i=m,$ implies that $A_{m+1}\in\Aset_{k-m}.$  


Next, let us apply Theorem \ref{th1sptr} for the special triplet $\{A_{m+1},B,C_{m+1}\},$   the maximal curve $\mu_{k-m},$ and $\Aset_{k-m}:=\Xset\cap\mu_{k-m},\ \Bset_{n-k+m}:=\Xset\setminus\mu_{k-m}.$
Below we present the first two statements obtained and, after each one, we draw the corresponding conclusions.
\begin{enumerate}[label*=\arabic*.]
\setcounter{enumi}{0}
\setlength{\itemsep}{0mm}
\item
The line $\ell_{m+1}:=\ell_{A_{m+1}C_{m+1}}$ is a component of $\mu_{k-m},$ in particular $B\notin\ell_{m+1}$ and $C_{m+1}\in \Aset_{k-m}.$ 
\end{enumerate}
From here we conclude that the line $\ell_{m+1}$ differs from each of the lines
$\ell_1, \ldots,\ell_m.$ Indeed, the latter lines are components of $\mu:=\mu_k$ and if, say, $\ell_{m+1}=\ell_1,$ then we will have that $\ell_1$ is a component of $\mu_{k-m}=\frac{1}{\ell_1\cdots\ell_m}\ \mu.$ Hence $\ell_1$ is a multiple component of a maximal curve $\mu,$ which is contradiction.

Note that this proves (i) for $m+1,$ without the ``Moreover" part, since all the remaining statements of (i) hold by virtue of the induction hypothesis.
\begin{enumerate}[label*=\arabic*.]
\setcounter{enumi}{1}
\setlength{\itemsep}{0mm}
\item
The set $\dot{\bar\ell}_{m+1}:=\{M\in \ell_{m+1}\cap  \Xset: \mu_{k-m-1}(M)
\neq 0\}$ 
besides the nodes $A_{m+1}$ and $C_{m+1},$ contains exactly $n-k+m$ nodes of $\Bset_{n-k+m+1},$ from which $n-k$ are the zeroes of $p^\star_{B,\Bset}$ on the line $\ell_{m+1},$  and $m$ are the following points of intersection: 
$$D_{im+1}:=\ell_i\cap\ell_{m+1}\ \ \forall\ 1\le i \le m.$$ 
\end{enumerate}
Let us mention that the statement on $n-k+m$ nodes follows from Theorem \ref{th1sptr}, (ii), and  the formula
\begin{equation*}p^\star_{B,\Bset_{n-k+m}}=p^\star_{B,\Bset}\prod_{i=1}^m\ell_i,
\end{equation*}
 which is a direct consequence of Theorem \ref{th1sptr}, (iv).

\noindent This proves (ii). Note that the statements on the lines $\dot{\bar\ell}_i,\ 1\le i\le m,$ hold by virtue of the induction hypothesis.

Note also that above $n-k$ zeroes and $m$ points of intersection together make 
$n-k+m$ points. Since this number is the least possible (see the proof of Theorem \ref{th1sptr}), we conclude that all the mentioned
points are distinct. From here we get the ``Moreover" part of the statement (i).  

Next, let us verify that the statements (iii) to (vi) are valid.

First, in view of Proposition \ref{maxcomp}, with $m=1,$ we obtain that the curve 

$$\mu_{k-m-1}:=\frac{1}{\ell_{m+1}}\ \mu_{k-m}=\frac{1}{\ell_1\cdots\ell_{m+1}}\ \mu_{k}$$ is a maximal curve of degree $k-m-1.$ This gives (iii).
 

 Then, the relation \eqref{maxmax} implies that $\Bset_{n-k+m+1}$ is an  $(n-k+m+1)$-correct set, for which thus $\ell_{m+1}$ is a maximal line. The remaining equalities in (iv) and (v) follow from the definitions of $\mu_{k-m-1}$ and ${\dot{\bar{\ell}}}_{m+1}.$  Hence, the statements (iv) and (v) are valid. 

Finally, let us verify (vi). On the line $\ell_{i}$ we have $n-k+m+1$ nodes from $\Bset_{n-k+m}$ which can be divided into the following three groups:

$G_1:$ the nodes $A_i, C_i,$ 

$G_2:$ the $n-k$ nodes which are the zeroes of $p^\star_{B,\Bset}$ on the line $\ell_i,$  and 

$G_3:$ the $m-1$ intersection points: $\{D_{ij},\ \forall\ 1\le j\le m,\ j\neq i\}.$

Note that, according to (ii), these three sets are disjoint. 

For the first two groups,  in view of (ii) and (iv), we have that

$G_1\cup G_2\subset \dot{\bar\ell}_i\subset\Bset_{n-k+i}\subset\Bset_{n-k+m}.$

For the third group, similarly we obtain

$G_3\subset\dot{\bar\ell}_i\cup\cdots\cup\dot{\bar\ell}_s\subset\Bset_{n-k+m}.$
This proves (vi).
\end{proof}

\section{Some results and corollaries for $2$-node lines}
Now, let us consider the case when the set $\Xset$ is $GC_n$ set and the special triplets arise from $2$-node lines.
\begin{theorem}\label{avel+}
Let $\Xset$ be $GC_n$ set. Let also $\mu:=\mu_k, \Bset:=\Bset_{n-k},\Aset:=\Aset_{k}, \Aset_i, \Bset_i$ be as in Theorem \ref{th1sptr}. 
Assume that the distinct special triplets $\{A_i,B,C_i\},$ $1\le i\le m,$ with a common node $B$ arise from $2$-node lines $\ell_{A_i,B},$ respectively, where $m\ge 2, A_i \in \Aset,\ B\in \Bset,$ and $C_i \in \Xset.$ 

Then the set $\Bset_{n-k+m}$ is $GC_{n-k+m}$ set. The nodes $B,C_1,\ldots, C_m$ are collinear and give one more maximal line of $\Bset_{n-k+m}: \ \ell_{BC_1}.$  Thus the set $\Bset_{n-k+m}$ has $m+1$ maximal lines.
\end{theorem}
\begin{proof}
It suffices to consider only the case $m=2:$

Suppose that $\Xset$ is a $GC_n$ set and the special triplets $\{A_i, B, C_i\}$ arise from the 2-node lines $\ell_{A_iB},\ i=1,2.$ Let us prove that the three nodes $B,\ C_1$ and $C_2$ are collinear. 

Consider the special case  
$k=n, \mu=\mu_n=p^\star_{B,\Xset},$
of Theorem. We have that  
$\Bset_{n-k+2}\equiv \Bset_2=\{B,A_1,C_1,A_2,C_2,D_{12}\}$ 
is 
$GC_2$ set. Let $\ell$ be a $3$-node (maximal) line the node $D_{12}$ uses in $\Bset_2.$ First, let us prove that $B\in \ell.$ Assume, by way of contradiction, that $B\notin \ell.$ Then, we get that three nodes from the set  $\{A_1,C_1,A_2,C_2\}$ are collinear. Therefore the three nodes contain either the couple $A_1,C_1$  or $A_2,C_2.$ In other words $\ell$ coincides either with the line $\ell_1$ or $\ell_2.$ This is a contradiction since, according to (ii),  the last two lines intersect at the node $D_{12},$ which differs from $A_1, A_2, C_1$ and $C_2.$

Thus we have that $B\in \ell.$ Since $\ell_{A_1B}$ and   $\ell_{A_2B}$ are $2$-node lines we conclude that $A_1,A_2\notin \ell.$ Therefore $C_1,C_2\in \ell.$
\end{proof}


It is worth mentioning the following
\begin{corollary} \label{cor2}
Let $\Xset$ be an $n$-correct set. Assume that $2$-node lines: $\ell_{A_1B}$ and $\ell_{A_2B}$ with a common node $B,$
are used by nodes $C_1$ and  $C_2$, respectively, where $A_i,B,C_i\in \Xset.$  

Then $C_1\neq C_2.$ Also, $C_1=A_2 \iff C_2= A_1.$

In  the case $C_1\neq A_2$ (and therefore $C_2\neq A_1$) the following hold: 
\begin{enumerate}
\setlength{\itemsep}{0mm}
\item
The lines $\ell_1:=\ell_{A_1C_1}$ and $\ell_2:=\ell_{A_2C_2}$ are different components of $p^\star_{B,\Xset}.$
\item
The curve
$$\mu_{n-2}:=\frac{1}{\ell_1 \ell_2}p^\star_{B,\Xset}$$
is a maximal curve of degree $n-2.$
\item
The set $\Bset_2:=\Xset\setminus \mu_{n-2}=\{A_1,A_2,B,C_1,C_2,D\}$ is a $2$-correct set, where $\ell_1, \ell_2$ are maximal lines and 
\begin{equation}\label{C=cap22}D:=D_{12}=\ell_1\cap \ell_2.\end{equation}
\item
If $\Xset$ is $GC_n$ set then the nodes $B,C_1,C_2$ are collinear and $\Bset_2$ is a $GC_2$ set. 
\end{enumerate}
\end{corollary}
\begin{figure}
     \centering
     \begin{subfigure}[b]{0.4\textwidth}
         \centering
         \includegraphics[width=\textwidth]{Figure4.png}
         \caption{if $\Xset$ is an $n$-correct set}
         \label{fig.3a}
     \end{subfigure}
     \hfill
     \begin{subfigure}[b]{0.4\textwidth}
         \centering
         \includegraphics[width=\textwidth]{Figure5.png}
         \caption{if $\Xset$ is a $GC_n$ set}
         \label{fig.3b}
     \end{subfigure}
            \caption{The set $\Bset_2$}
        \label{Fig4}
\end{figure}

In Fig. \ref{Fig4} we have the set $\Bset_2$ if $\Xset$ is (a): an $n$-correct and (b): a $GC_n$ set. Note that the set $\Bset_2$ in the case (a) is not $GC_2$ set, since (only) the fundamental polynomial of $D_{12}$ is not a product of linear factors. While  the set $\Bset_2$ in the case (b) is a $GC_2$ set. 


Next we consider the problem whether a set of distinct used $2$-node lines corresponds to a set of distinct special triplets.
\begin{proposition} \label{n>2}
Let $\Xset$ be $GC_n$ set.
Let distinct $2$-node lines $\ell_{A_i,B},\  1\le i\le m,$  be used by the nodes $C_i,$ where $A_i,B,C_i\in\Xset.$ The following hold:
\begin{enumerate}
\setlength{\itemsep}{0mm}
\item
If $m=2$ then $C_1=A_2,$ which is equivalent to $C_2=A_1,$ implies that the corresponding  two triplets coincide.
While $C_1\neq A_2$ implies that the two triplets are distinct.
\item
If $m\ge 3$ then the corresponding special triplets $\{A_i,B,C_i\},\  1\le i\le m$ are distinct too.
\end{enumerate}
\end{proposition}
\begin{proof}
(i) Assume that $m=2.$ Consider the triplets $\{A_1,B,C_1\}$ and $\{A_2,B,C_2\}.$
Since $A_1\neq A_2$ therefore, in view of Proposition \ref{sptr2}, we have that 
\begin{equation} \label{abbc}\{A_1,B,C_1\}=\{A_2,B,C_2\} \iff  A_1=C_2.\end{equation}
(ii) Now assume that $m=3.$ Consider three distinct used $2$-node lines: $\ell_{A_1B}, \ell_{A_2B}, \ell_{A_3B},$ with a common node $B,$ where $A_1,A_2,A_3,B\in \Xset.$
Assume, that these lines are
used by the nodes $C_1, C_2$ and  $C_3$, respectively, where $C_i\in \Xset.$

Now, we readily get that the nodes  $C_1,C_2$ and $C_3$ are distinct  since a node, in view of Corollary \ref{2nodeuse}, cannot use two different $2$-node lines. 

Then, let us prove that the nodes  $C_1,C_2,$ and $C_3$ are different from the nodes  $A_1,A_2,$ and $A_3.$ Of course we have that $C_i\neq A_i,\ i=1,2,3,$ since the node $C_i$ uses the line  $\ell_{A_iB}.$

Next, assume by way of contradiction that, say, $A_1=C_2.$ Then, in view of \eqref{abbc},  we obtain that $A_2=C_1.$ Now, we readily get that $C_3$ is different from $A_1$ and $A_2.$ 

Indeed, assume conversely that, say, $C_3=A_2.$ Then we have $C_3=C_1,$ which is contradiction.

Since $C_3\neq A_2$ (and therefore $C_2\neq A_3$, from \ref{abbc}),  in view of  Theorem \ref{avel+},  we get that the nodes $B, C_2, C_3$ are collinear.

Therefore we have that $C_3\in\ell_{C_2B}\equiv\ell_{A_1B}.$ Hence, the line $\ell_{A_1B}$ is not a $2$-node line, which is contradiction.


Thus the nodes  $C_1, C_2$ and $C_3$ are different from the nodes  $A_1,A_2,$ and $A_3$ and therefore the special triplets $\{A_i,B,C_i\},\  i=1,2,3,$ are distinct.

Now let us consider the case of general $m\ge 3.$ Here, by considering different triplets of nodes, we get readily that $C_1,\ldots C_m$ are distinct and
they are different from the nodes  $A_1,\ldots,A_m.$ Consequently, the corresponding triplets  $\{A_i,B,C_i\},\  1\le i\le m,$ are distinct.
\end{proof}
In Fig. \ref{Fig6} we have two used $2$-node lines $\ell_{A_iB},\ i=1,2,$ with a common node, in a $3$-correct set, which correspond to the same special triplet $\{A_1,B,A_2\}.$ Indeed, $A_2$ uses $\ell_{A_1B}$ and $A_1$ uses $\ell_{A_2B},$ since $p^\star_{A_2}=\ell_{A_1B}\ell_1\ell_2$ and $p^\star_{A_1}=\ell_{A_2B}\ell_1\ell_2.$ While, again in $3$-correct set,  in each Fig. \ref{Fig4} a) and b) we have two used $2$-node lines $\ell_{A_iB},\ i=1,2,$ with a common node, which correspond to two distinct special triplets $\{A_i,B,C_i\},\ i=1,2,$ respectively.
\begin{corollary}
\emph{(i)} Let $\Xset$ be an $n$-correct set, $n\ge 2.$ Then there are at most $n$ special triplets with a fixed common node $B\in \Xset.$ In the case there are $n$ such triplets we have exactly $n$ maximal lines in $\Xset$ not passing through $B.$

\emph{(ii)} Let $\Xset$ be $GC_n$ set, $n\ge 3.$ Then there are at most $n$ used $2$-node lines with a fixed common node $B\in\Xset.$ In the case there are $n$ such lines we have exactly $n+1$ maximal lines in $\Xset,$ from which $n$ not passing through $B$ and a maximal line passing through $B.$ Thus the set $\Xset$ is Carnicer-Gasca set.
\end{corollary}
\begin{figure}
\begin{center}
\includegraphics[width=6.0cm,height=3.0cm]{Figure6.png}
\end{center}
\caption{Two used $2$-node lines correspond to one special triplet} \label{Fig6}
\end{figure}

\begin{proof}
(i) Let us apply Theorem \ref{prp:m2ket}, (i), with the maximal curve $\mu=p^\star_{B,\Xset}$ of degree $n.$ Then for the number $m$ of distinct special triplets with a fixed common node $B$ we have the inequality $m\le n.$

In the case there are 
$n$ such triplets we obtain from Theorem \ref{prp:m2ket}, (vi), that there are $n$ maximal lines: $\ell_1,\ldots,\ell_n$ in the $n$-correct set $\Bset_n\equiv \Xset$ not passing through $B.$ There is no another maximal line not passing through $B,$ which, in view of Proposition \ref{prp:n+1ell}, would contradict the existence of the fundamental polynomial $p^\star_B.$ 

(ii) In this case we obtain from Proposition \ref{n>2} that the triplets corresponding to $2$-node lines are distinct. Then Theorem \ref{avel+} implies that there are $n+1$ maximal lines in $\Xset,$ described in Corollary (ii). From here we readily obtain that there are exactly $n+1$ maximal lines in $\Xset.$ Indeed, otherwise if there are $n+2$ such lines then $\Xset$ becomes Chung-Yao set, where the only used lines are the maximal lines and therefore no $2$-node line is used. 
\end{proof} 
In Fig. \ref{Fig1} we have Carnicer-Gasca set ($GC_3$ set) and three used $2$-node lines with a common node. We have also exactly $3$ maximal lines: $\ell_i,\ i=1,2,3,$ not passing through $B$ and a maximal line: $\ell_0,$ passing through $B.$

At the end let us formulate a result concerning $GC_6$ sets.
\begin{corollary}\label{6655} Let $\Xset$ be $GC_6$ set without a maximal line. 
Suppose that there are three used $2$-node lines: $\ell_i:=\ell_{A_iB},\ i=1,2,3,$ with a common node $B,$ where $A_i,B\in \Xset.$ Then there are three $6$-node lines, not intersecting at nodes of $\Xset.$
\end{corollary}
\begin{proof} Assume, by way of contradiction, that there are the mentioned three used $2$-node lines: $\ell_i,\ i=1,2,3.$
Then, according to Proposition \ref{n>2}, (iii),  the three special triplets corresponding to the three used $2$-node lines are distinct. Thus, by using Theorem \ref{prp:m2ket}, we obtain that the curve 
$$\mu_{3}:=\frac{1}{\ell_1 \ell_2 \ell _3}p^\star_{B,\Xset}$$
is a maximal curve of degree $3.$ Hence $\mu_3$ passes through $d(6,3)=18$ nodes and is a product of three distinct lines. Since these lines are not maximal, we readily get that each passes through $6$ nodes and non of these nodes can be common for two or three of them.
\end{proof}




 \begin{thebibliography}{99}

\bibitem{dB07}
 {C. de Boor, Multivariate polynomial interpolation: conjectures concerning GC-sets. Numer. Algorithms
{\bf 45} 113-125 (2007).}

\bibitem{CG01}
J. M. Carnicer and M. Gasca, Planar Configurations with Simple lagrange formula. In: Lyche, T., Schumaker,
L.L. (eds.) Mathematical Methods in CAGD: Oslo 2000, pp. 55–62. Vanderbilt University
Press, Nashville (2001).

\bibitem{CG03}
J. M. Carnicer and M. Gasca,
{On Chung and Yao's geometric characterization for bivariate polynomial
interpolation}. In: Lyche, T., Mazure, M.-L., Schumaker, L.L. (eds.) Curve and surface design: Saint
Malo 2002, pp. 21-30. Nashboro Press, Brentwood (2003).


\bibitem{CY77}
K.~C.~Chung and T.~H.~Yao,
{On lattices admitting unique Lagrange interpolations},
SIAM J.\ Numer.\ Anal.\ {\bf 14} 735--743 (1977).


\bibitem{HJZ09b}
H. Hakopian, K. Jetter, and G. Zimmermann,
A new proof of the
                Gasca-Maeztu conjecture for $n=4.$ Journal of Approximation
                Theory {\bf 159}(2) 224-242 (2009).

\bibitem{HakTor}
H. Hakopian and S. Toroyan, On the uniqueness of algebraic curves passing through  $n$-independent nodes. \textit{New York J. Math.}, {\bf 22} 441--452 (2016).

\bibitem{HV}
{H. Hakopian and N. Vardanyan,
On the basic properties of $GC_n$ sets,
Journal of Knot Theory and Its Ramifications,
{\bf 29}  1-26 (2020).}

\bibitem{HVV}
H. Hakopian G. Vardanyan and N. Vardanyan,
On maximal curves of $n$-correct sets, arXiv:2507.11207 [math.NA], (2025)
https://doi.org/10.48550/arXiv.2507.11207

\bibitem{Raf}
L. Rafayelyan,
{Poised nodes set constructions on algebraic curves},
East J.\ Approx.\ {\bf 17} 285--298 (2011).

\end{thebibliography}

%%%%%%%%%%%%%%%%%%%%%%%%%%%%%%%%%%%%%%%%%%%%%%%%%%%%%%%%%%%%%%%%%%%%%%%%
%%%%%%%%%%%%%%%%%%%%%%%%%%%%%%%%%%%%%%%%%%%%%%%%%%%%%%%%%%%%%%%%%%%%%%%%%%
%%%%%%%%%%%%%%%%%%%%%%%%%%%%%%%%%%%%%%%%%%%%%%%%%%%%%%%%%%%%%%%%%%%%%%%%%%%%%%
%%%%%%%%%%%%%%%%%%%%%%%%%%%%%%%%%%%%%%%%%%%%%%%%%%%%%%%%%%%%%%%%%%%%%%%%%%%%%%

\end{document}





















\begin{lemma}\label{lemma01}
The line $3^*$ can be used by at most four nodes in $\Aset.$  
If it is used by exactly four nodes Then, all they belong to the set 
$\alpha_4\setminus 3^*.$

Moreover, the following two statements hold:

(i) At most one node from each line $\alpha_{i},\ 1\le i\le 3,$ can use the line $3^*.$

(ii) If a node from $\alpha_4\setminus 3^*$ uses $3^*$ Then, all four nodes from it use $3^*.$
  \end{lemma}

\begin{proof} Note that $3^*$ is a $6$-node line hence, according to Proposition \ref{prp:6810}, it can be used by at most seven nodes of $X.$ On the other side it is a maximal line for $\Bset$ and hence is used by three nodes from the set $\Bset.$ Thus we get that the line $3^*$ can be used by at most four nodes in $\Aset.$  

Next, let us prove the statement $(ii).$ Suppose that a node from the set $\alpha_4\setminus 3^*$ does not use $3^*$ and assume by way of contradiction that a node from  $\alpha_4\setminus 3^*$ uses $3^*.$

Consider the set $\Nset_{3^*}.$ We have that $\Bset\cap \Nset_{3^*}=\emptyset$ and \begin{equation}\label{aa*}0<\#\left[\alpha_4\cap \Nset_{3^*}\right]\le 3.\end{equation}

Now, we easily get a $5$-fundamental polynomial for the set $\Nset_{3^*}.$ Namely, let $A\in \alpha_4\cap \Nset_{3^*},$ we take $p^\star_A=\alpha_1 \alpha_2 \alpha_3 \ell_1\ell_2,$ where the product $\ell_1\ell_2,$ does not pass through $A$ and passes through all the nodes of the set $\alpha_4\cap \Nset_{3^*}\setminus \{A\}$ (whose cardinality does not exceed $2$). 

This contradicts to Theorem ...

Now, to complete the proof it is enough to prove the statement $(i).$ Suppose that two nodes  from the set $\alpha_i,$ where $1\le i\le 3$ use the line $3^*.$ Then, there are five nodes in the set  $\alpha_i\setminus 3^*$ thus at least one of them uses $3^*.$ Therefore we get the relaion \eqref{aa*}, where $4$ is replaced by $i.$ This, in the same way as above, finishes the proof.
\end{proof}
Then, a $2$-node line $\ell_0:=\ell_{AB},$ where $A\in\Aset,B\in\Bset,$ is used by a node $C\in \Xset$ if and only if the the following four conditions hold.
Suppose that $\Xset$ is $GC_6$-set without a maximal line Then, we have the following.

(a) The four line components of $\mu_4$ contain $6,6,5,5$ primary nodes respectively. Each of the six nodes of $\Sset$ has the m-d sequence $(6,6,5,5,3,2).$ In the respective line sequence the first four lines are the line components of $\mu_4.$ The last two lines are the lines the node of $\Sset$ uses inside the $GC_2$-set $\Sset.$

(b) Either $\Sset=\Bset$ or $\Sset\subset \Aset.$ In the latter case we have a double $GC_2$ construction.

















In the sequel we will need the following results:
\begin{theorem}[case m=1:  \cite{HakTor}, Thm.~4.2; case m=2: \cite{HK}, Thm.~3] \label{th:-1+1}
Let $m=1$ or $2.$ Assume that ${\mathcal X}$ is an $n$-independent set of $d(n, k-m)+m$ nodes with $1+m\le k\le n-1.$ Then, at most $2m$ different curves of degree
$\le k$ pass through all the nodes of ${\mathcal X}.$

Moreover, there are such $2m$ curves for the set ${\mathcal X}$ if and only if all the nodes of ${\mathcal X}$ but $m$ lie in a maximal curve of degree $k-m.$
\end{theorem}


\begin{theorem}[\cite{HKV}, Thm.~2.5, \cite{H}, Thm.~3.2]\label{th:-2+3}
Assume that ${\mathcal X}$ is an $n$-independent set of $d(n, k-2)+3$ nodes, $3\le k\le n-1.$ Then, at most $3$ linearly independent curves of degree
$\le k$ pass through all the nodes of ${\mathcal X}.$

 Moreover, there are such three curves for the set ${\mathcal X}$ if and only if all the nodes of ${\mathcal X}$  lie in a  curve of degree $k-1,$ or all the nodes of  ${\mathcal X}$ but three lie in a (maximal) curve of degree $k-2.$
\end{theorem}
Below we bring a characterization of $n$-dependent sets $\Xset$ with $\#\Xset\le 3n.$
\begin{theorem}[\cite{HM}, Thm.~5.1]\label{th:3n}
A set $\Xset$ consisting of at most $3\,n$ nodes is $n$-dependent if
and only if one of the following conditions holds.
\vspace{-2mm}
\begin{enumerate}
\setlength{\itemsep}{0mm}
\item
$n+2$ nodes are collinear,
\item
$2n+2$ nodes belong to a (possibly reducible) conic,
\item
$\#\Xset = 3n,$ and there exist $\gamma\in\Pi_3$ and $\sigma\in\Pi_n$
such that $\Xset = \gamma \cap \sigma$\,.
\end{enumerate}
\end{theorem}

\begin{corollary}\label{cor:3n-1}
A set $\Xset$ consisting of at most $3n-1$ nodes is $n$-dependent if
and only if either $n+2$ nodes are collinear, or
$2n+2$ nodes belong to a (possibly reducible) conic.
\end{corollary}


Consider special $n$-correct sets: $GC_n$ sets,  defined by Chung and Yao:
\begin{definition}[\cite{CY77}]
An n-correct set $\Xset$ is called $GC_n$ \emph{set}, if the $n$-fundamental polynomial of each node $A\in\Xset$ is a product of
$n$ linear factors.
\end{definition}

Now, we are in a position to present the Gasca-Maeztu, or briefly GM
\begin{conjecture}[\cite{GM82}, 1982] \label{GM}
Any $GC_n$ set contains $n+1$ collinear nodes.
\end{conjecture}

So far, the GM conjecture has been confirmed to be true only for $n \leq 5$. The case $n=2$ is trivial. The case $n=3$ was established by M. Gasca and J. I. Maeztu in \cite{GM82}. The case
$n = 4$ was proved by J. R.  Busch \cite{B90}. Other proofs of this case have been published since Then, (see e.g. \cite{CG01}, \cite{HJZ09b}). The case $n=5$ was proved by H. Hakopian, K. Jetter and G. Zimmermann \cite{HJZ14}. Recently G. Vardanyan provided a simpler and shorter proof for this case \cite{GV2}.

In this paper we make a step in proving the Gasca-Maeztu conjecture for $n=6$ (see Proposition \ref{thm:main3}). The analogue of this step was crucial in the proof of the case  $n=5$ (see \cite{HJZ14}, Prop. 3.12; \cite{GV2}, Prop. 2.8).


\begin{definition} Let $\Xset$ be an $n$-correct set.
We say, that a node $A\in\Xset$
\emph{uses a line} $\ell$, if  $p_{A}^\star=\ell q,\quad q\in\Pi_{n-1}.$
\end{definition}

Since the fundamental polynomial in an $n$-correct set is unique we
get
\begin{lemma} \label{lm}
Suppose $\Xset$ is an $n$-correct set and a node  $A\in \Xset$ uses a line $\ell.$ Then, $\ell$ passes through at least two nodes from $\Xset$, at which $q$ from the above definition does not vanish.
 \end{lemma}


\begin{definition} \label{def:Nll}
For a given set of lines $\ell_1,\dots,\ell_k,$ we define
$\Nset_{\ell_1,\dots,\ell_k}$
to be the set of those nodes in $\Xset$ which do not lie in any of the
lines $\ell_i$, and for which at least one of the lines is not used.
\end{definition}
\noindent In the case of one line $\ell$ we have
$$ \Nset_\ell = \left\{ A\in\Xset : A \notin \ell, \text{ and $A$ is
  not using $\ell$}\right\}.$$
\begin{proposition}[\cite{HJZ09b}, Thm. 3.2]\label{prp:kell}
Assume that $\Xset$ is a $GC_n$ set, and
$\ell_1,\dots,\ell_k$ are lines. Then, the following hold for $\Nset =
\Nset_{\ell_1,\dots,\ell_k}$\,.
\vspace{-2mm}
\begin{enumerate}
\setlength{\itemsep}{0mm}
\item
If $\, \Nset$ is nonempty, Then, it is essentially   $(n-k)$-dependent.
\item
$\Nset = \emptyset$ if and only if the product $\ell_1\cdots\ell_k$ is a maximal curve.
\end{enumerate}
\end{proposition}

\noindent For $k=1$ this result has been proved by Carnicer and Gasca \cite{CG01}.

Assume that $\Xset_i$ is a set of $k_i$ collinear points:
$$\Xset_i\subset \ell_i,\quad \#\Xset=k_i,\ i=1,2,3,\quad \ell_i \hbox{ is a line}.$$
Assume also that non of the points is an intersection point of the lines.

 
\subsection{The main result}\label{case665532}

In ... we proved the following
\begin{proposition}\label{thm:main3} Assume that $\Xset$ is a $GC_6$ set with no maximal line. Then, for no node in $\Xset$ the m-d sequence is $(6,6,6,4,3,2)$.
\end{proposition}

In this paper we will prove the following
\begin{proposition}\label{thm:main3} Assume that $\Xset$ is a $GC_6$ set with no maximal line. Then, for no node in $\Xset$ the m-d sequence is $(6,6,5,5,3,2)$.
\end{proposition}

Assume by way of contradiction that for a node
in $\Xset$  the m-d sequence is  $(6,6,5,5,3,2)$. Let $(\alpha_1,\ldots,\alpha_6)$ be a respective m-line sequence.

Set  $\Xset=\Aset \cup \Bset$ (see Fig. \ref{case1}) with
\begin{equation*}\label{6}
  \Aset = \Xset \cap \{\alpha_1 \cup \alpha_2 \cup \alpha_3 \cup \alpha_4\}, \quad
  \Aset_{12} = \Xset \cap \{\alpha_1 \cup \alpha_2\}, \quad
   \Aset_{34} = \Xset \setminus  \Aset_{12},\\   
   \end{equation*}
   \begin{equation*}
    \#\Aset = 22, \quad \#\Aset_{12} = 12, \quad\#\Aset_{34} = 10,\quad\text{and} \quad
  \Bset = \Xset \setminus \Aset, \quad  \#\Bset = 6.
\end{equation*}

Denote $\Lset_4:=\{\alpha_1,\alpha_2,\alpha_3,\alpha_4\}.$
Note that no intersection point of the three lines of $\Lset_3$ belongs to $\Xset$.
The following is the analogue of  \cite{HJZ14}, Lemma 3.2.
\begin{lemma}\label{lem:Bset}
\nix\vspace{-2mm}
\begin{enumerate}
\setlength{\itemsep}{0mm}
\item
The set $\Bset$ is a $GC_2$ set, and each node $B\in\Bset$ uses the four lines of $\Lset_4$ and the two lines it uses within $\Bset$, i.e.,
\begin{equation}\label{ban}
  p_{B,\Xset}^\star = \alpha_1\alpha_2\alpha_3 \alpha_4p_{B,\Bset}^\star\,.
\end{equation}

\item
Only one node  in $\Aset$ may use the line $\alpha_1.$  This latter node can belong only to  the line $\alpha_2.$\\ 
Likewise, only one node in $\Aset$ may use the line $\alpha_2.$ This latter node can belong only to  the line $\alpha_1.$ 
No node in $\Aset$ uses the line $\alpha_3$ or $\alpha_4.$

\item
If the lines $\alpha_1$ and $\alpha_2$ intersect Then, the intersection point is not a node of $\Xset.$ Likewise, if the lines $\alpha_3$ and $\alpha_4$ intersect Then, the intersection point is not a node of $\Xset.$ 

\item 
If the lines $\alpha_3$ and $\alpha_4$ intersect the line $\alpha_1$ (or the line $\alpha_2$) at nodes of $\Xset$, Then, these nodes coincide. In other words, the two lines intersect the third line at the same node of $\Xset.$ 
\end{enumerate}
\end{lemma}

\begin{proof}
(i) Here we use the relation \eqref{maxmax} and the fact that $\mu_4:=\alpha_1\alpha_2\alpha_3 \alpha_4$ is a maximal curve.

(ii) Without loss of generality assume conversely that $A \in \alpha_3$ uses the line $\alpha_1.$ Then, $p_{A}^\star = \alpha_1\,q,$ where $q \in \Pi_5.$ It is easily seen that $q$ has (6,5,4) primary zeros in the lines $(\alpha_2 , \alpha_4, \alpha_3).$ Therefore, in view of Corollary \ref{cor22}, we obtain that $p_{A}^\star =\alpha_1\alpha_2\alpha_3\alpha_4r,\ r \in \Pi_2,$ which is a contradiction.


Next, assume that two nodes $A_1,A_2 \in \alpha_2$ use the line $\alpha_1.$ 
In view of Lemma \ref{lem:mm-1} for a point $T_0\in\alpha_3\setminus \Xset$ there is a polynomial
$$p_0\in linear span \{p^\star_{A_1},p^\star_{A_2}\},\quad p_0\neq 0,$$

\noindent such that $p_0(T_0)=0.$

Now, it is easily seen that $p_0$ has (6,5,4) primary zeros in the lines $(\alpha_3 , \alpha_4,\alpha_2).$ Therefore, in view of Corollary \ref{cor22}, we obtain that $p_{A}^\star =\alpha_1\alpha_2\alpha_3\alpha_4r,\ r \in \Pi_2,$ which is a contradiction.

Then,  assume that $A \in \alpha_4$ uses the line $\alpha_3.$ Then, $p_{A}^\star = \alpha_3\,q,$ where $q \in \Pi_5.$ It is easily seen that $q$ has (6,5,4) primary zeros in the lines $(\alpha_1 , \alpha_2,\alpha_4).$ Therefore, in view of Corollary \ref{cor22}, we obtain that $p_{A}^\star =\alpha_1\alpha_2\alpha_3\alpha_4r,\ r \in \Pi_3,$ which is a contradiction.

(iii) Assume by way of contradiction that the lines $\alpha_1$ and $\alpha_2$ intersect at a node $A\in\Xset.$ Then, $A$ is a primary node for the line $\alpha_1,$ while it is a secondary node and is not counted for the line $\alpha_1.$ Therefore the line $\alpha_2$ actually contains $7$ nodes which is a contradiction.

Then, assume by way of contradiction that the lines $\alpha_3$ and $\alpha_4$ intersect at a node $A\in\Xset.$ Then, $A$ is a primary node for the line $\alpha_3,$ while it is a secondary node and is not counted for the line $\alpha_4.$ Therefore the line $\alpha_4$ actually contains $6$ nodes. Now, if we interchange $\alpha_3$ and $\alpha_4$ in the line sequence Then, we will get greater distribution sequence for $6,6,6,4,3,2,$ which is a contradiction.

(iv) Assume by way of contradiction that the lines $\alpha_3$ and $\alpha_4$ intersect
the line $\alpha_1$ at the different nodes $A_1,A_2\in\Xset.$ 
Therefore, as above, we get that the lines $\alpha_3$ and $\alpha_4$ actually contains $6$ nodes. 
 Now, if we reorder the first four lines the line sequence as follows $\alpha_3, \alpha_4,\alpha_2,\alpha_1$ Then, we will get greater distribution sequence $6,6,6,4,3,2,$ which is a contradiction.
\end{proof}



\item
Suppose that $\Xset$ is $GC_n$ set and $\mu:=\mu_{n-2}$ is a maximal curve of degree $n-2$ and  hence the set $\Bset:=\Xset\setminus \mu$ is a $GC_2$ set. 
Assume that no $2$-node line $\ell_0:=\ell_{AB},$ where $A\in\Aset,\ B\in\Bset,$ is used by a node $C\in \Aset.$ 
Assume also that two $2$-node lines: $\ell_{FA_1}$ and $\ell_{FA_2}$ with a common node $F,$
are used by the nodes $B_1$ and  $B_2$, respectively, where $F,A_i,B_i\in \Xset.$ Finally, assume that $B_1\neq A_2$ (therefore $B_2\neq A_1$) and $\Sset$ is the $GC_2$ set from the item (i).  

Then, we have that either $\Sset=\Bset$ or $\Sset\subset \Aset.$







First, let us verify that $B\in \Aset,$ i.e., $\mu(B)=0.$ Indeed, conversely assume that $B\in \Bset.$ Then we have that 
\begin{equation}\label{1102}p^\star_{B,\Xset}=p^\star_{B,\Bset} \mu.
\end{equation}
The node $B$ uses the line $\ell_{FA},$ hence the latter divides either $p^\star_{B,\Bset}$ or $\mu.$ But neither of these two can take place since the line $\ell_{FA}$ passes through exactly one node of $\Bset.$
\begin{proof}
Since the triplets are distinct, Proposition \ref{sptr2} readily implies that $$\{A_1,C_1\}\cap \{A_2,C_2\}=\emptyset.$$
We have that $A_1,B,C_1$ is a special triplet. Therefore, in view of Theorem \ref{th1sptr} we obtain that:
 
\begin{enumerate}[label*=\arabic*.]
\item
The line $\ell_1:=\ell_{A_1C_1},$ is a component of $\mu,$ in particular $C_1\in \Aset.$ 
\item
The curve $\mu_{k-1}:=\frac{1}{\ell_1}\ \mu$ is a maximal curve of degree $k-1.$ 
\item
The set $\dot{\bar\ell}_1:=\{M\in\Xset\cap {\ell_1} : \mu_{k-1}(M)\neq 0\}$ besides the nodes $A_1$ and $C_1,$ contains exactly $n-k$ more nodes of $\Xset,$ 
which coincide with the $n-k$ zeroes of $p^\star_{F,\Bset}$ on the line $\ell_1.$ 
\end{enumerate}

Next, we apply Theorem \ref{th1sptr} for the special triplet $A_2,B,C_2,$   the maximal curve $\mu_{k-1},$ and $\Aset_{k-1}=\Xset\cap\mu_{k-1},\ \Bset_{n-k+1}=\Xset\setminus\mu_{k-1}.$

In view of  $A_2\neq A_1, C_1$ and Theorem \ref{th1sptr}, (vii), we conclude that $A_2\notin \dot{\bar\ell}_1,$ hence $\mu_{k-1}(A_2)=0,$ i.e., $A_2\in\Aset_{k-1}:=\Xset\cap\mu_{k-1}.$  

Thus we obtain the following:
\begin{enumerate}[label*=\arabic*.]
\setcounter{enumi}{3}
\setlength{\itemsep}{0mm}
\item
The line $\ell_2:=\ell_{A_2C_2}$ is a component of $\mu_{k-1},$ in particular $C_2\in \Aset_{k-1}.$ 
\item
The curve $\mu_{k-2}:=\frac{1}{\ell_2}\ \mu_{k-1}$ is a maximal curve of degree $k-2.$ 
\item
The set $\dot{\bar\ell}_2:=\{M\in\Xset\cap {\ell_2} : \mu_{k-2}(M)\neq 0\}$ besides the nodes $A_2$ and $B_2,$ contains exactly $n-k+1$ more nodes of $\Xset.$ 
\item
These $n-k+1$ nodes coincide with the $n-k$ zeroes of $p^\star_{B,\Bset_{n-k}}$ on the line $\ell_2$ and the point of intersection of the lines $\ell_1$ and $\ell_2.$
\end{enumerate}
The last statement follows from the formula
$$p^\star_{F,\Bset_{n-k+1}}=\ell_1p^\star_{F,\Bset_{n-k}}.$$
Here we used Theorem Theorem \ref{th1sptr}, (ii), as well as (v), according to which the line $\ell_1$ is a maximal line for $\Bset_{n-k+1}.$

Now, from the statements 1 and 4 we get (i), from 2 and 5 we get (ii), from 3, 6, and 7, we get (iii). Note that the line $\ell_1$ obtains one more node in $\Bset_{n-k+2},$ namely the intersection point $\ell_1\cap\ell_2.$ Thus each of the lines $\ell_1$ and $\ell_2$ contains exactly $n-k+3$ nodes in $\Bset_{n-k+2}.$  Thus (iv) and (v) are proved.
 
Next, to prove (vi) consider the special case  
$k=n, \mu=\mu_n=p^\star_{B,\Xset},$
of Theorem. We have that  
$\Bset_{n-k+2}\equiv \Bset_2=\{B,A_1,C_1,A_2,C_2,D\}$ 
is 
$GC_2$ set. Let $\ell$ be a $3$-node (maximal) line the node $D$ uses. *** First, let us prove that $B\in \ell.$ Assume, by way of contradiction, that $B\notin \ell.$ Then, we get that three nodes from the set  $\{A_1,C_1,A_2,C_2\}$ are collinear. Therefore the three nodes contain either $A_1,C_1$  or $A_2,C_2.$ In other words $\ell$ coincides either with the line $\ell_1$ or $\ell_2.$ This is a contradiction since, according to \eqref{D12}, the last two lines pass through the node $D.$

Thus we have that $B\in \ell.$ Since $\ell_{A_1B}$ and   $\ell_{A_2B}$ are $2$-node lines we conclude that $A_1,A_2\notin \ell.$ Therefore $C_1,C_2\in \ell.$
\end{proof}


